% templates/attractor.tex.mustache
% AUTO-GENERATED Attractor Report Section
% This file is rendered by Mustache template from trajectory data
% DO NOT EDIT MANUALLY

\section{Attractor Dynamics and Low-Dimensional Phase Space}

\subsection{Overview}

The boundary layer is analyzed as a low-dimensional dynamical system through projection onto dominant singular vectors. The following section presents the reconstructed phase space trajectory for campaign \texttt{ ALL } using the v0.1 baseline.

\begin{description}
    \item[Campaign:] ALL
    \item[Baseline Version:] v0.1
    \item[Total Samples:] 6682
    \item[Mean Entropy:] $H_{\text{mean}} = 0.223$
\end{description}

\subsection{3D Phase Space Projection}

Figure~\ref{fig:attractor_ALL} visualizes the three-dimensional trajectory $(\eta_1, \eta_2, \eta_3)$. To highlight the true multi-regime orbital evolution, timesteps are connected sequentially by a continuous historical path line. The trajectory is colored using a high-contrast colormap driven by the \textbf{Attractor Spin} diagnostic ($\Omega(t)$), explicitly isolating the directionality of regime transitions.

\begin{figure}[htbp]
    \centering
    \begin{tikzpicture}
        \begin{axis}[
            width=0.9\textwidth,
            height=0.55\textwidth,
            view={45}{24},
            grid=major,
            xlabel={$\eta_1$},
            ylabel={$\eta_2$},
            zlabel={$\eta_3$},
            title={Low-Rank Attractor Trajectory (\texttt{ ALL })},
            colormap/jet,
            colorbar,
            point meta=explicit,
            colorbar style={title={$\Omega(t)$}, title style={yshift=5pt}},
            tick label style={font=\small},
            label style={font=\small},
            title style={font=\small}
        ]
            % 1. DRAW THE ORBIT LINES: Connects data sequentially to reveal the phase trajectory loops
            \addplot3[
                mesh,
                line width=0.8pt,
                mark=none,
                scatter,
                scatter src=explicit,
                shader=interp
            ] table[
                col sep=comma,
                x=eta_1,
                y=eta_2,
                z=eta_3,
                meta=attractor_spin
            ] {../../data/outputs/regime_scatterplots_all.csv};

        \end{axis}
    \end{tikzpicture}
    \caption{3D attractor state orbit trajectory. Color variation tracks Attractor Spin ($\Omega$), separating shear breakout acceleration (cool blue) from radiative collapse configurations (warm red).}
    \label{fig:attractor_ALL}
\end{figure}

\subsection{Physical Coordinate Interpretation}

This phase-space view acts as the compressed state engine of the atmospheric boundary layer (ABL), projecting high-dimensional tower profiles into three dominant, orthogonal coordinates:

\begin{itemize}
    \item $\eta_1$ (\textbf{Bulk Background Mode}): Captures the mean inversion strength, nocturnal cooling depth, and overall boundary-layer depth evolution.
    \item $\eta_2$ (\textbf{Shear \& LLJ Mode}): Captures shear production, alongside low-level jet (LLJ) speed and height modulation.
    \item $\eta_3$ (\textbf{Vertical Structural Curvature Mode}): Captures submesoscale gravity-wave excitation and the precursor structural signatures of intermittent turbulence events.
\end{itemize}

\subsubsection{Geometric Reading Guide}
\begin{itemize}
    \item \textbf{Tight closed loops} indicate near-repeatable diurnal cycling (convective daytime mixing to stable nighttime recovery and back).
    \item \textbf{Broad, tangled trajectories} indicate non-stationary forcing, intermittency, or mesoscale disruption.
    \item \textbf{Color evolution by $\Omega(t)$} isolates transition directionality: positive spin values (warm tones) indicate stable collapse paths; negative values (cool tones) mark turbulent regeneration breakouts.
\end{itemize}

\subsection{Trajectory Metrics}

The mean low-rank coefficients across all samples are:
\begin{align}
    \langle \eta_1 \rangle &= 0.574 \\
    \langle \eta_2 \rangle &= 0.284 \\
    \langle \eta_3 \rangle &= -0.077
\end{align}

The magnitude of the attractor vector $\|\boldsymbol{\eta}\|$ evolves continuously throughout the campaign, reflecting the transient structural adjustments driven by diurnal heating and nocturnal radiative cooling cycles.

Operationally, peaks in $\|\boldsymbol{\eta}\|$ often coincide with rapid boundary-layer reconfiguration, while spin excursions identify windows where similarity-based single-regime assumptions are least reliable.

\subsection{Information-Theoretic Manifold Complexity Evaluation}

To preserve algorithmic invariance across divergent energy scales, rank truncation is enforced with a machine-precision floor criterion ($s_i \ge \epsilon_{\text{mach}} s_1$). This execution yields:

\begin{itemize}
    \item \textbf{Constrained modes:} $N_c = 3$
    \item \textbf{Unconstrained nullspace modes:} $N_0 = 0$
    \item \textbf{Condition number:} $\kappa = 6.85$
    \item \textbf{Shannon effective dimension:} $D_{\text{eff}} = 2.0537$
\end{itemize}

The entropy-derived $D_{\text{eff}} = 2.0537$ quantifies how many modes are actively carrying structure in practice, rather than only counting nominal basis size. Larger values indicate broader modal participation and lower compressibility of the observed boundary-layer state.

\label{sec:attractor_ALL}